% from ega4-2-fr.tex
\oldpage[IV-2]{29}
Les propriétés (v) et (vi) ne sont mises que pour mémoire, étant conséquences
de (1.1.4), (1.1.6) et (2.3.7). En ce qui concerne les autres, la condition
est nécessaire en vertu de (2.4.3). Inversement, supposons par exemple que
$f'$ soit universellement ouvert, et soit $Y_1\to Y$ un morphisme
quelconque~; posons $X_1=X\times_Y Y_1$ et $f_1=f_{(Y_1)}$. D'autre part,
posons
\[
Y'_1=Y_1\times_Y Y',\qquad
X'_1=X'\times_{Y'}Y'_1=X_1\times_{Y_1}Y'_1,
\qquad f'_1=f'_{(Y'_1)}=f_{1(Y'_1)}:X'_1\to Y'_1~;
\]
comme $g'=g_{(Y')}:Y'_1\to Y'$ est fidèlement plat et quasi-compact
(2.2.13 et 1.1.2), et que $f'$ est universellement ouvert, $f'_1$ est ouvert
et il résulte de (2.6.2) que $f_1$ est ouvert~; donc $f$ est universellement
ouvert. Même raisonnement pour les autres cas.

On notera encore ici que l'on peut remplacer «~fidèlement plat~» par
«~quasi-fidèlement plat~» et, lorsque $g$ est en outre localement de
présentation finie, ou que $g$ est surjectif et $S$ discret, on peut remplacer
la propriété (iii) par~:
\begin{enumerate}
\item[(iii \emph{bis})] universellement bicontinu.
\end{enumerate}

\subsection{Permanence de diverses propriétés des morphismes par descente
fidèlement plate}
\label{subsection:IV.2.7-fr}

\begin{proposition}[2.7.1]
\label{IV.2.7.1-fr}
Soient $f:X\to Y$ un $S$-morphisme de $S$-préschémas, $g:S'\to S$ un
morphisme fidèlement plat et quasi-compact, $X'=X_{(S')}$, $Y'=Y_{(S')}$,
$f'=f_{(S')}:X'\to Y'$. Considérons, pour un morphisme, la propriété d'être~:
\begin{enumerate}
\item[(i)] séparé~;
\item[(ii)] quasi-séparé~;
\item[(iii)] localement de type fini~;
\item[(iv)] localement de présentation finie~;
\item[(v)] de type fini~;
\item[(vi)] de présentation finie~;
\item[(vii)] propre~;
\item[(viii)] un isomorphisme~;
\item[(ix)] un monomorphisme~;
\item[(x)] une immersion ouverte~;
\item[(xi)] une immersion quasi-compacte~;
\item[(xii)] une immersion fermée~;
\item[(xiii)] affine~;
\item[(xiv)] quasi-affine~;
\item[(xv)] fini~;
\item[(xvi)] quasi-fini~;
\item[(xvii)] entier.
\end{enumerate}
Alors, si $P$ désigne une des propriétés précédentes, pour que $f$ possède la
propriété $P$, il faut et il suffit que $f'$ la possède.
\end{proposition}

Il a été démontré aux chapitres I\textsuperscript{er}, II et dans le chapitre
IV, \S~1, que si $f$ possède une des propriétés $P$ ci-dessus, il en est de
même de $f'$ (sans hypothèse sur le morphisme $g:S'\to S$). Il reste donc à
démontrer la réciproque~; comme la projection $Y'\to Y$ est

